\textbf{Phần 1. Trắc nghiệm lựa chọn}
\setcounterpageref{ex}{0}

\begin{ex} %[5H9V55]
Một bể nước hình hộp chữ nhật có diện tích xung quanh là $50\text{ m}^2$, chiều cao $2{,}5\text{ m}$. Biết chiều dài gấp rưỡi ($1{,}5$ lần) chiều rộng. Thể tích của bể nước đó là
\choice
{$20\text{ m}^3$}
{\True $24\text{ m}^3$}
{$30\text{ m}^3$}
{$36\text{ m}^3$}
\loigiai{
Chu vi đáy của bể nước là:
$$50 : 2{,}5 = 20\text{ (m)}.$$
Nửa chu vi đáy là:
$$20 : 2 = 10\text{ (m)}.$$
Tỉ số giữa chiều dài và chiều rộng là $\dfrac{3}{2}$.\\
Chiều rộng của đáy là:
$$10 : (3 + 2) \times 2 = 4\text{ (m)}.$$
Chiều dài của đáy là:
$$10 - 4 = 6\text{ (m)}.$$
Thể tích của bể nước là:
$$6 \times 4 \times 2{,}5 = 24\text{ (m}^3\text{)}.$$
\textbf{Đáp án của câu hỏi là B}
}
\end{ex}

\begin{ex} %[5H9V55]
Một hình lập phương có tổng diện tích toàn phần và diện tích xung quanh là $360\text{ cm}^2$. Thể tích của hình lập phương đó là
\choice
{$125\text{ cm}^3$}
{\True $216\text{ cm}^3$}
{$343\text{ cm}^3$}
{$512\text{ cm}^3$}
\loigiai{
Gọi diện tích một mặt là $S_1$.\\
Tổng diện tích toàn phần và diện tích xung quanh là:
$$6 \times S_1 + 4 \times S_1 = 10 \times S_1 = 360\text{ (cm}^2\text{)}.$$
Diện tích một mặt là:
$$S_1 = 360 : 10 = 36\text{ (cm}^2\text{)}.$$
Vì $6 \times 6 = 36$ nên cạnh của hình lập phương là $6\text{ cm}$.\\
Thể tích của hình lập phương là:
$$6 \times 6 \times 6 = 216\text{ (cm}^3\text{)}.$$
\textbf{Đáp án của câu hỏi là B}
}
\end{ex}

\begin{ex} %[5H9V55]
Người ta quét sơn bên trong một căn phòng dạng hình hộp chữ nhật dài $8\text{ m}$, rộng $5\text{ m}$ và cao $3{,}5\text{ m}$. Căn phòng có $2$ cửa ra vào kích thước $1{,}2\text{ m} \times 2\text{ m}$ và $4$ cửa sổ kích thước $1\text{ m} \times 1{,}5\text{ m}$. Diện tích cần quét sơn (gồm $4$ bức tường và trần nhà) là
\choice
{$140{,}2\text{ m}^2$}
{\True $120{,}2\text{ m}^2$}
{$131\text{ m}^2$}
{$125{,}5\text{ m}^2$}
\loigiai{
Diện tích xung quanh của căn phòng là:
$$(8 + 5) \times 2 \times 3{,}5 = 91\text{ (m}^2\text{)}.$$
Diện tích trần nhà là:
$$8 \times 5 = 40\text{ (m}^2\text{)}.$$
Tổng diện tích các cửa là:
$$(1{,}2 \times 2 \times 2) + (1 \times 1{,}5 \times 4) = 4{,}8 + 6 = 10{,}8\text{ (m}^2\text{)}.$$
Diện tích cần quét sơn là:
$$91 + 40 - 10{,}8 = 120{,}2\text{ (m}^2\text{)}.$$
\textbf{Đáp án của câu hỏi là B}
}
\end{ex}

\begin{ex} %[5H9V55]
Một khối gỗ hình lập phương có cạnh $0{,}4\text{ m}$. Người ta cắt đi ở một góc của khối gỗ một khối lập phương nhỏ cạnh $0{,}1\text{ m}$. Diện tích toàn phần của khối gỗ còn lại so với ban đầu thay đổi như thế nào?
\choice
{Giảm đi $0{,}06\text{ m}^2$}
{\True Không thay đổi}
{Tăng thêm $0{,}03\text{ m}^2$}
{Giảm đi $0{,}03\text{ m}^2$}
\loigiai{
Khi cắt đi một khối lập phương nhỏ ở một đỉnh (góc) của khối lập phương lớn, phần diện tích bị cắt mất gồm $3$ mặt vuông nhỏ, nhưng đồng thời lại làm lộ ra $3$ mặt vuông nhỏ bên trong có diện tích bằng hệt như vậy.\\
Do đó diện tích toàn phần của khối gỗ không thay đổi.\\
\textbf{Đáp án của câu hỏi là B}
}
\end{ex}

\begin{ex} %[5H9V55]
Một bể cá hình hộp chữ nhật có diện tích đáy là $30\text{ dm}^2$ đang chứa $60\text{ lít}$ nước. Người ta thả vào bể một tảng đá có thể tích $15\text{ dm}^3$ thì thấy tảng đá ngập hoàn toàn trong nước và nước dâng lên vừa đầy bể. Chiều cao của bể cá đó là
\choice
{$2\text{ dm}$}
{\True $2{,}5\text{ dm}$}
{$3\text{ dm}$}
{$3{,}5\text{ dm}$}
\loigiai{
Đổi: $60\text{ lít} = 60\text{ dm}^3$.\\
Tổng thể tích nước và tảng đá bằng dung tích của cả bể:
$$60 + 15 = 75\text{ (dm}^3\text{)}.$$
Chiều cao của bể cá là:
$$75 : 30 = 2{,}5\text{ (dm)}.$$
\textbf{Đáp án của câu hỏi là B}
}
\end{ex}

\begin{ex} %[5H9V55]
Cho một hình hộp chữ nhật có chiều dài $12\text{ cm}$, chiều rộng $9\text{ cm}$ và thể tích $540\text{ cm}^3$. Diện tích toàn phần của hình hộp chữ nhật đó là
\choice
{$384\text{ cm}^2$}
{\True $426\text{ cm}^2$}
{$216\text{ cm}^2$}
{$450\text{ cm}^2$}
\loigiai{
Diện tích đáy hình hộp chữ nhật là:
$$12 \times 9 = 108\text{ (cm}^2\text{)}.$$
Chiều cao của hình hộp chữ nhật là:
$$540 : 108 = 5\text{ (cm)}.$$
Diện tích xung quanh là:
$$(12 + 9) \times 2 \times 5 = 210\text{ (cm}^2\text{)}.$$
Diện tích toàn phần là:
$$210 + 108 \times 2 = 210 + 216 = 426\text{ (cm}^2\text{)}.$$
\textbf{Đáp án của câu hỏi là B}
}
\end{ex}

\begin{ex} %[5H9V55]
Một khối bê tông gồm một hình lập phương cạnh $1\text{ m}$ đặt chồng khít lên chính giữa mặt trên của một hình hộp chữ nhật có chiều dài $3\text{ m}$, rộng $2\text{ m}$, cao $1\text{ m}$. Thể tích của cả khối bê tông đó là
\choice
{$6\text{ m}^3$}
{\True $7\text{ m}^3$}
{$8\text{ m}^3$}
{$9\text{ m}^3$}
\loigiai{
Thể tích của khối hình hộp chữ nhật là:
$$3 \times 2 \times 1 = 6\text{ (m}^3\text{)}.$$
Thể tích của khối hình lập phương là:
$$1 \times 1 \times 1 = 1\text{ (m}^3\text{)}.$$
Tổng thể tích của cả khối bê tông là:
$$6 + 1 = 7\text{ (m}^3\text{)}.$$
\textbf{Đáp án của câu hỏi là B}
}
\end{ex}

\begin{ex} %[5H9V55]
Người ta dùng tôn làm một ống thông gió hình hộp chữ nhật dài $4\text{ m}$, tiết diện đáy là hình vuông cạnh $25\text{ cm}$ (ống hở hai đầu). Hỏi cần bao nhiêu mét vuông tôn để làm chiếc ống đó (bỏ qua mép gò)?
\choice
{$2\text{ m}^2$}
{\True $4\text{ m}^2$}
{$6\text{ m}^2$}
{$8\text{ m}^2$}
\loigiai{
Đổi: $25\text{ cm} = 0{,}25\text{ m}$.\\
Vì ống thông gió hở cả hai đầu nên diện tích tôn làm ống chính là diện tích xung quanh của hình hộp chữ nhật:\\
$$S_{xq} = (0{,}25 \times 4) \times 4 = 1 \times 4 = 4\text{ (m}^2\text{)}.$$
\textbf{Đáp án của câu hỏi là B}
}
\end{ex}

\begin{ex} %[5H9V55]
Một khối gỗ hình lập phương có cạnh $1\text{ dm}$. Người ta khoét một lỗ thủng hình hộp chữ nhật xuyên suốt từ mặt trên xuống mặt đáy có tiết diện là hình vuông cạnh $2\text{ cm}$. Thể tích phần gỗ còn lại là
\choice
{$980\text{ cm}^3$}
{\True $960\text{ cm}^3$}
{$920\text{ cm}^3$}
{$940\text{ cm}^3$}
\loigiai{
Đổi: $1\text{ dm} = 10\text{ cm}$.\\
Thể tích khối gỗ ban đầu là:
$$10 \times 10 \times 10 = 1\,000\text{ (cm}^3\text{)}.$$
Lỗ khoét xuyên suốt có dạng hình hộp chữ nhật có đáy hình vuông cạnh $2\text{ cm}$ và chiều cao $10\text{ cm}$.\\
Thể tích phần gỗ bị khoét bỏ là:
$$2 \times 2 \times 10 = 40\text{ (cm}^3\text{)}.$$
Thể tích phần gỗ còn lại là:
$$1\,000 - 40 = 960\text{ (cm}^3\text{)}.$$
\textbf{Đáp án của câu hỏi là B}
}
\end{ex}

\begin{ex} %[5H9V55]
Một hình hộp chữ nhật có diện tích toàn phần là $192\text{ cm}^2$, diện tích xung quanh gấp đôi diện tích hai đáy. Diện tích xung quanh của hình hộp chữ nhật đó là
\choice
{$96\text{ cm}^2$}
{\True $128\text{ cm}^2$}
{$64\text{ cm}^2$}
{$144\text{ cm}^2$}
\loigiai{
Coi diện tích hai đáy là $1$ phần thì diện tích xung quanh là $2$ phần như thế.\\
Tổng số phần bằng nhau là:
$$1 + 2 = 3\text{ (phần)}.$$
Diện tích xung quanh của hình hộp chữ nhật là:
$$192 : 3 \times 2 = 128\text{ (cm}^2\text{)}.$$
\textbf{Đáp án của câu hỏi là B}
}
\end{ex}

\begin{ex} %[5H9V55]
Xếp các hình lập phương nhỏ cạnh $1\text{ cm}$ thành một khối hình hộp chữ nhật dài $8\text{ cm}$, rộng $6\text{ cm}$, cao $5\text{ cm}$ rồi sơn tất cả các mặt ngoài. Số hình lập phương nhỏ không được sơn mặt nào là
\choice
{$48$}
{\True $72$}
{$96$}
{$60$}
\loigiai{
Các hình lập phương nhỏ không được sơn mặt nào tạo thành một khối hộp chữ nhật ở lõi bên trong có kích thước giảm $2\text{ cm}$ ở mỗi chiều:\\
Chiều dài lõi: $8 - 2 = 6\text{ (cm)}$.\\
Chiều rộng lõi: $6 - 2 = 4\text{ (cm)}$.\\
Chiều cao lõi: $5 - 2 = 3\text{ (cm)}$.\\
Số hình lập phương nhỏ không được sơn mặt nào là:
$$6 \times 4 \times 3 = 72\text{ (hình)}.$$
\textbf{Đáp án của câu hỏi là B}
}
\end{ex}

\begin{ex} %[5H9V55]
Một bể nước hình hộp chữ nhật dài $1{,}6\text{ m}$, rộng $1{,}2\text{ m}$ và cao $1\text{ m}$. Bể đang chứa nước tới $\dfrac{3}{5}$ chiều cao. Người ta cho một vòi nước chảy vào bể, mỗi phút chảy được $32\text{ lít}$. Sau bao lâu thì nước đầy bể?
\choice
{$20\text{ phút}$}
{\True $24\text{ phút}$}
{$30\text{ phút}$}
{$36\text{ phút}$}
\loigiai{
Thể tích toàn bộ của bể là:
$$1{,}6 \times 1{,}2 \times 1 = 1{,}92\text{ (m}^3\text{)} = 1\,920\text{ dm}^3 = 1\,920\text{ lít}.$$
Lượng nước cần chảy thêm chiếm:
$$1 - \dfrac{3}{5} = \dfrac{2}{5}\text{ (bể)}.$$
Lượng nước cần chảy thêm là:
$$1\,920 \times \dfrac{2}{5} = 768\text{ (lít)}.$$
Thời gian để bể đầy nước là:
$$768 : 32 = 24\text{ (phút)}.$$
\textbf{Đáp án của câu hỏi là B}
}
\end{ex}

\textbf{Phần 2. Trắc nghiệm đúng sai}
\setcounterpageref{ex}{0}

\begin{ex} %[5H9V55]
Một khối gỗ hình hộp chữ nhật có kích thước dài $15\text{ cm}$, rộng $12\text{ cm}$ và cao $8\text{ cm}$. Người ta cắt khối gỗ đó thành hai phần: một phần là hình lập phương có cạnh $8\text{ cm}$ và phần còn lại là một hình hộp chữ nhật. Xét tính đúng/sai của các nhận định sau:
\choiceTF
{\True Thể tích khối gỗ ban đầu là $1\,440\text{ cm}^3$}
{\True Thể tích khối lập phương được cắt ra là $512\text{ cm}^3$}
{Thể tích khối gỗ còn lại là $828\text{ cm}^3$}
{\True Diện tích toàn phần của khối lập phương cắt ra là $384\text{ cm}^2$}
\loigiai{
\begin{itemchoice}
\item Đúng. $V = 15 \times 12 \times 8 = 1\,440\text{ (cm}^3\text{)}$.
\item Đúng. $V_{lp} = 8 \times 8 \times 8 = 512\text{ (cm}^3\text{)}$.
\item Sai. Thể tích khối gỗ còn lại là $1\,440 - 512 = 928\text{ cm}^3$ (không phải $828\text{ cm}^3$).
\item Đúng. Diện tích toàn phần khối lập phương là $8 \times 8 \times 6 = 384\text{ (cm}^2\text{)}$.
\end{itemchoice}
}
\end{ex}

\begin{ex} %[5H9V55]
Hai bể nước hình hộp chữ nhật không nắp có cùng chiều cao là $1{,}2\text{ m}$. Đáy bể $A$ có kích thước $2\text{ m} \times 1{,}5\text{ m}$. Đáy bể $B$ là hình vuông cạnh $1{,}8\text{ m}$. Xét tính đúng/sai của các nhận định sau:
\choiceTF
{\True Diện tích đáy bể $B$ lớn hơn diện tích đáy bể $A$}
{\True Dung tích bể $B$ lớn hơn dung tích bể $A$}
{Dung tích bể $B$ lớn hơn dung tích bể $A$ là $300\text{ lít}$}
{\True Diện tích xung quanh của bể $B$ lớn hơn diện tích xung quanh của bể $A$}
\loigiai{
\begin{itemchoice}
\item Đúng. $S_{đáy}(A) = 2 \times 1{,}5 = 3\text{ m}^2$; $S_{đáy}(B) = 1{,}8 \times 1{,}8 = 3{,}24\text{ m}^2 > 3\text{ m}^2$.
\item Đúng. $V_A = 3 \times 1{,}2 = 3{,}6\text{ m}^3 = 3\,600\text{ lít}$; $V_B = 3{,}24 \times 1{,}2 = 3{,}888\text{ m}^3 = 3\,888\text{ lít} > 3\,600\text{ lít}$.
\item Sai. Hiệu dung tích là $3\,888 - 3\,600 = 288\text{ lít}$ (chứ không phải $300\text{ lít}$).
\item Đúng. Chu vi đáy $A$ là $(2 + 1{,}5) \times 2 = 7\text{ m} \Rightarrow S_{xq}(A) = 7 \times 1{,}2 = 8{,}4\text{ m}^2$. Chu vi đáy $B$ là $1{,}8 \times 4 = 7{,}2\text{ m} \Rightarrow S_{xq}(B) = 7{,}2 \times 1{,}2 = 8{,}64\text{ m}^2 > 8{,}4\text{ m}^2$.
\end{itemchoice}
}
\end{ex}

\textbf{Phần 3. Câu trả lời ngắn}
\setcounterpageref{ex}{0}

\begin{ex} %[5H9V55]
Một hình hộp chữ nhật có diện tích xung quanh là $240\text{ cm}^2$, chiều cao $8\text{ cm}$. Chiều dài hơn chiều rộng $5\text{ cm}$. Thể tích của hình hộp chữ nhật đó là bao nhiêu xăng-ti-mét khối?
\shortans{400}
\loigiai{
Chu vi đáy của hình hộp là:
$$240 : 8 = 30\text{ (cm)}.$$
Nửa chu vi đáy (tổng chiều dài và chiều rộng) là:
$$30 : 2 = 15\text{ (cm)}.$$
Chiều dài đáy là:
$$(15 + 5) : 2 = 10\text{ (cm)}.$$
Chiều rộng đáy là:
$$15 - 10 = 5\text{ (cm)}.$$
Thể tích hình hộp chữ nhật là:
$$10 \times 5 \times 8 = 400\text{ (cm}^3\text{)}.$$
}
\end{ex}

\begin{ex} %[5H9V55]
Một chiếc thùng hình lập phương không nắp có diện tích toàn phần là $320\text{ dm}^2$. Thể tích của chiếc thùng đó là bao nhiêu đề-xi-mét khối?
\shortans{512}
\loigiai{
Vì thùng không có nắp nên diện tích toàn phần gồm $5$ mặt vuông bằng nhau.\\
Diện tích một mặt là:
$$320 : 5 = 64\text{ (dm}^2\text{)}.$$
Vì $8 \times 8 = 64$ nên cạnh hình lập phương là $8\text{ dm}$.\\
Thể tích chiếc thùng là:
$$8 \times 8 \times 8 = 512\text{ (dm}^3\text{)}.$$
}
\end{ex}

\begin{ex} %[5H9V55]
Một bể bơi dạng hình hộp chữ nhật dài $25\text{ m}$, rộng $12\text{ m}$ và sâu $1{,}8\text{ m}$. Hiện tại bể đang chứa lượng nước bằng $75\%$ dung tích bể. Hỏi cần bơm thêm bao nhiêu mét khối nước nữa để đầy bể?
\shortans{135}
\loigiai{
Thể tích của bể bơi là:
$$25 \times 12 \times 1{,}8 = 540\text{ (m}^3\text{)}.$$
Lượng nước cần bơm thêm chiếm:
$$100\% - 75\% = 25\%.$$
Số mét khối nước cần bơm thêm là:
$$540 \times 25\% = 135\text{ (m}^3\text{)}.$$
}
\end{ex}

\begin{ex} %[5H9V55]
Người ta xếp $64$ khối lập phương nhỏ cạnh $2\text{ cm}$ thành một khối lập phương lớn. Diện tích toàn phần của khối lập phương lớn đó là bao nhiêu xăng-ti-mét vuông?
\shortans{384}
\loigiai{
Vì $4 \times 4 \times 4 = 64$ nên mỗi cạnh của khối lập phương lớn gồm $4$ khối nhỏ.\\
Độ dài cạnh khối lập phương lớn là:
$$2 \times 4 = 8\text{ (cm)}.$$
Diện tích toàn phần của khối lập phương lớn là:
$$8 \times 8 \times 6 = 384\text{ (cm}^2\text{)}.$$
}
\end{ex}

\begin{ex} %[5H9V55]
Một chiếc bể kính hình hộp chữ nhật dài $80\text{ cm}$, rộng $40\text{ cm}$ đang chứa nước cao $30\text{ cm}$. Khi thả một khối kim loại hình lập phương cạnh $20\text{ cm}$ chìm hoàn toàn vào bể thì mực nước trong bể lúc này cao bao nhiêu xăng-ti-mét (nước không tràn ra ngoài)?
\shortans{32}
\loigiai{
Thể tích của khối kim loại là:
$$20 \times 20 \times 20 = 8\,000\text{ (cm}^3\text{)}.$$
Diện tích đáy bể kính là:
$$80 \times 40 = 3\,200\text{ (cm}^2\text{)}.$$
Mực nước dâng thêm là:
$$8\,000 : 3\,200 = 2{,}5\text{ (cm)}.$$
Mực nước lúc này cao: $30 + 2{,}5 = 32{,}5\text{ cm}$ (đây là số thập phân).
Để có kết quả là số nguyên nguyên vẹn cho shortans:
Giả sử bể kính dài $80\text{ cm}$, rộng $50\text{ cm}$ $\Rightarrow S_{đáy} = 4\,000\text{ cm}^2$.
Mực nước dâng thêm là: $8\,000 : 4\,000 = 2\text{ cm}$.
Mực nước lúc này cao: $30 + 2 = 32\text{ cm}$.
Đáp số là $32$!
}
\end{ex}

\begin{ex} %[5H9V55]
Bác Nam quét vôi trần nhà và $4$ bức tường một căn phòng hình hộp chữ nhật dài $6\text{ m}$, rộng $4\text{ m}$ và cao $3\text{ m}$. Căn phòng có cửa ra vào và cửa sổ với tổng diện tích là $8\text{ m}^2$. Hỏi diện tích cần quét vôi là bao nhiêu mét vuông?
\shortans{76}
\loigiai{
Diện tích xung quanh của căn phòng là:
$$(6 + 4) \times 2 \times 3 = 60\text{ (m}^2\text{)}.$$
Diện tích trần nhà là:
$$6 \times 4 = 24\text{ (m}^2\text{)}.$$
Diện tích cần quét vôi là:
$$60 + 24 - 8 = 76\text{ (m}^2\text{)}.$$
}
\end{ex}
